% Original analytic teaching diagram, 2026-11-12.
% A random variable X: Omega -> R. The preimage X^{-1}(B) of a Borel set B (here two pieces in Omega)
% must be an event; its probability is the mass the distribution P_X puts on B.
\definecolor{VB}{HTML}{7188A3}\definecolor{VBF}{HTML}{E5EAF0}
\definecolor{VG}{HTML}{829C8B}\definecolor{VGF}{HTML}{E8EEE9}
\definecolor{VY}{HTML}{B4A16F}\definecolor{VYF}{HTML}{F1ECDD}
\definecolor{VR}{HTML}{A7787D}\definecolor{VRF}{HTML}{F0E5E6}
\definecolor{VP}{HTML}{9588A5}\definecolor{VPF}{HTML}{ECE8F0}
\definecolor{VI}{HTML}{3B4252}\definecolor{VM}{HTML}{6B7280}\definecolor{VL}{HTML}{4C4D4F}
\begin{tikzpicture}[x=1mm,y=1mm,>=Stealth,font=\figurefont\fontsize{8}{10}\selectfont,
lab/.style={text=VM,font=\figurefont\fontsize{7.5}{9}\selectfont}]
% sample space
\draw[VL,line width=.7pt] (18,17) ellipse (17 and 14);
\node[text=VI] at (4,29) {$\Omega$};
\fill[VBF] (10,22) ellipse (5 and 3.6); \draw[VB,line width=.7pt] (10,22) ellipse (5 and 3.6);
\fill[VBF] (24,9) ellipse (6 and 3.4); \draw[VB,line width=.7pt] (24,9) ellipse (6 and 3.4);
\node[text=VI] at (24,20) {$X^{-1}(B)$};
\draw[VL,line width=.4pt] (19.5,18.6) -- (13,21.5);
\draw[VL,line width=.4pt] (24,18.4) -- (24,11.5);
\fill[VI] (10,22) circle (.6); \fill[VI] (25,8.5) circle (.6); \fill[VI] (30,20) circle (.6); \fill[VI] (8,11) circle (.6);
% value space
\draw[->,VL,line width=.6pt] (48,4) -- (104,4);
\node[text=VI,anchor=west] at (104.5,4) {$\mathbb{R}$};
\draw[VB,line width=2.2pt] (72,4) -- (86,4);
\node[text=VI] at (79,7.5) {$B$};
\fill[VI] (79,4) circle (.6); \fill[VI] (75,4) circle (.6); \fill[VI] (58,4) circle (.6); \fill[VI] (95,4) circle (.6);
% arrows of the map
\draw[->,VL,line width=.5pt] (10.6,22.4) to[bend left=18] (74.6,4.8);
\draw[->,VL,line width=.5pt] (25.8,9.6) to[bend left=24] (78.6,4.9);
\draw[->,VM,line width=.45pt,densely dashed] (30.6,20.2) to[bend left=14] (94.6,4.8);
\draw[->,VM,line width=.45pt,densely dashed] (8.6,10.6) to[bend right=22] (57.6,3.4);
\node[text=VI] at (52,24) {$X$};
\node[text=VI] at (79,-3) {$P_X(B)=\mathbb{P}\bigl(X^{-1}(B)\bigr)$};
\end{tikzpicture}
